\section{Problem formulation}
\label{sec:problem}

\subsection{Flat-folded states}
\label{sec:problem-states}

A flat-folded state is an isometric map of the sheet into the plane together with a layer
ordering. The map is fixed by the crease pattern; the ordering is not. A single crease pattern
therefore admits many folded states that are geometrically identical and distinct as embeddings,
and determining the ordering is NP hard even when a valid mountain-valley assignment is given
\citep{bern-hayes}.

This distinction is the reason the task cannot be reduced to geometry. Over a folding sequence the
crease graph's topology never changes: the same vertices, edges and faces exist at step one and at
step nineteen. What changes is everything else. Vertex coordinates are reflected at every step.
The overlap structure, which faces lie above which, changes at every step and follows from the
geometry. The layer ordering changes at every step and does \emph{not} follow from the geometry;
it is chosen, and it is where the decisions live. The number of layers accumulates, which is the
formal content of ``it gets harder as it goes''.

\subsection{The action space}
\label{sec:problem-actions}

Simple folds come in named variants, and any benchmark in this area has to say which it uses,
because the complexity results differ between them.

\begin{table}[h]
\caption{Named variants of the simple fold \citep{arkin-map}.}
\label{tab:action-space}
\begin{center}
\begin{tabular}{ll}
\multicolumn{1}{c}{\bf MODEL}  &\multicolumn{1}{c}{\bf DEFINITION} \\
\hline \\
One-layer simple fold    & Rotate a single layer $\pm 180^\circ$ about a line \\
Some-layers simple fold  & Rotate a contiguous run of layers \\
All-layers simple fold   & Rotate every layer the line crosses, as sheet metal bends \\
Finite vs infinite line  & Whether the fold line is a full line or a segment \\
\end{tabular}
\end{center}
\end{table}

CP2Seq instantiates the all-layers and some-layers tiers. The all-layers tier is the simpler
object and it buys three properties at once, the third of which is rarely named. Because the whole
stack moves as a rigid body, no two connected pieces of paper ever move relative to each other
except at the fold line itself, where paper is allowed to bend. Self-intersection is impossible,
the layer ordering is forced rather than chosen, and \textbf{tearing is impossible by
construction.}

Moving only some layers brings all three back, and tearing becomes the binding constraint. If a
moving face is joined to a stationary face along a crease, and that crease is not on the fold
line, the fold would rip the sheet. This is decidable only in the coordinates of the original
square, because adjacency is a fact about the original sheet that survives any amount of folding
and cannot be recovered from the current positions of the polygons alone.

An example makes the tier difference concrete, and it needs only two folds. Fold the right half of
a square over to the left across the vertical midline. The result is a two-layer stack whose top
layer is joined to the bottom layer along the entire crease at the midline. Now select the top
layer alone. Folding its upper half down across the horizontal midline would tear the sheet,
because the moving piece is joined to stationary paper along a vertical crease segment that is
perpendicular to the fold line. Folding its left quarter across a vertical line does not tear,
because the moving piece's only join to stationary paper lies on the fold line itself. Same state,
same layer selection, opposite verdicts, and the only difference is the orientation of the line.

------this 3 paragraphs are over emphasizing on whether all layers fold or some layer'fold  it
gonna tear or not. is that the only point in this part????

\subsection{The task}
\label{sec:problem-task}

Given a crease pattern, produce a sequence of simple folds that reproduces it. Generating an
instance costs one forward pass of the engine; solving it is NP hard. That asymmetry is the
paper's foundation.

---------

It is tempting
to say that legal moves are rare and that this is what makes the task hard. Enumeration refutes
it: legal partial folds \emph{grow} with depth, from 22 at two layers to 332 at thirty-eight, six
times the 56 all-layers folds available at the same state. What falls is the hit rate of uniform
random proposal, from 11.9\% to 3.8\% by seventeen layers, because the space being sampled grows
faster than the legal set inside it. That is a fact about a sampler, not about origami, and
conflating the two would put a false statement about branching factor into the paper.

---------this paragraph should be listed properly and give a proper graph to illustrate it!!!!!!!!

